\chapter*{CONCLUSION AND RECOMMENDATIONS}
\phantomsection
\addcontentsline{toc}{chapter}{CONCLUSION AND RECOMMENDATIONS}

\section*{Conclusion}

This secondary analysis of Viet Nam MICS6 found that HPV vaccine awareness among women aged 15--29 years was moderate, but self-reported any-dose uptake was low. Overall, 62.4\% of women had heard of HPV vaccination, while only 7.5\% reported having received any HPV vaccine dose. Among women who were aware of the vaccine, 88.0\% remained unvaccinated.

Awareness and any-dose uptake were concentrated among wealthier women, whereas the awareness--vaccination gap was concentrated among poorer aware women. These findings indicate that barriers after awareness---including affordability, availability, distance, and service delivery constraints---are likely to be important, although these mechanisms were not directly measured in MICS6.

The results provide a pre-program equity baseline rather than causal estimates. They show that public financing and delivery systems should be designed explicitly to reduce socioeconomic, geographic, educational, and ethnic inequities as HPV vaccination is expanded under the national EPI roadmap \cite{govvn2022nq104,baochinhphu2025epi}.

\section*{Recommendations}

\begin{enumerate}
    \item \textbf{Financing.} Public financing or strong subsidies should prioritize poor households, remote communes, and ethnic minority communities so that HPV vaccination does not remain primarily a private-market service for better-off families.
    \item \textbf{Delivery.} School-based delivery should be used where girls remain in school, while community sessions, mobile outreach, commune health stations, and adolescent or reproductive health services should reach out-of-school girls and young women missed by school-based platforms.
    \item \textbf{Communication.} Communication should address both awareness and practical access. Materials should use clear language, local languages where needed, and trusted channels such as commune health workers, village health collaborators, women's unions, schools, and local community leaders.
    \item \textbf{Equity monitoring.} Monitoring should report any-dose uptake and, when available, dose completion by wealth, residence, ethnicity, education, school attendance, province/region, and implementation area rather than national averages alone.
    \item \textbf{Data systems.} Future surveys and administrative systems should collect dose number, vaccination date, vaccine product, provider type, payment source, and reasons for non-vaccination to distinguish affordability, availability, acceptance, and completion barriers.
\end{enumerate}

\section*{Recommended methodological extensions}

The current thesis provides a valid national cross-sectional baseline, but several extensions could strengthen future manuscript versions or a journal submission.

\begin{enumerate}
    \item \textbf{Use multilevel survey models if geographic identifiers become available.} A hierarchical model with women nested within provinces or regions could estimate how much inequality is attributable to household-level socioeconomic position versus contextual health-system access. This is particularly relevant because future EPI implementation will be phased geographically.
    \item \textbf{Separate the awareness and uptake stages formally.} A sequential framework could model awareness first and vaccination conditional on awareness second, or use a hurdle-style approach. This would avoid interpreting low vaccination as a single process when the results show both information barriers and post-awareness access barriers.
    \item \textbf{Add interaction and effect-modification analyses.} Pre-specified interactions between wealth and residence, wealth and education, or ethnicity and residence could identify groups for targeted outreach, provided cell sizes remain adequate under survey weighting.
    \item \textbf{Improve causal clarity with a directed acyclic graph.} Future analyses should explicitly distinguish confounders, mediators, and equity stratifiers. For example, education and wealth may confound some access associations but may also lie on broader structural pathways; over-adjustment should be avoided when estimating total social inequality.
    \item \textbf{Use multiple imputation or missing-data sensitivity analysis if item missingness increases in expanded models.} The current complete-case adjusted models were acceptable because the analytic variables had limited missingness, but imputation would be preferable if adding additional reproductive or service-access variables.
    \item \textbf{Incorporate dose completion and service-channel variables when available.} The strongest limitation of MICS6 for HPV vaccination is the absence of detailed dose number, vaccine type, payment source, service location, and reasons for non-vaccination. These variables would allow separation of affordability, availability, acceptance, and completion barriers.
    \item \textbf{Plan post-introduction monitoring.} Once 2026--2028 EPI HPV activities produce administrative or survey data, repeat the same concentration-index and decomposition framework to test whether public financing reduces the pro-rich gradient observed in this pre-program baseline.
\end{enumerate}

\clearpage
